\chapter{Number Systems}

\section{Digital Representations}

\begin{definition}[Tuple]
    In mathematics, a tuple is a finite ordered sequence of elements. An $n$-tuple is a sequence of $n$ elements, where $n$ is a non-negative integer.
\end{definition}
In digital representations, a number is represented by an ordered $n$-tuple, where each element of the tuple is called a digit. The $n$-tuple is often referred to as a digit vector (or string of digits) and $n$ is called the presicison of the representation.

\subsection{Representation of Nonnegative Integers}
\begin{definition}[Digit-vector]
    A digit-vector is an $n$-tuple of digits representing the integer $x$ denoted as:
    \[
        X = (X_{n-1}, X_{n-2}, \ldots, X_0)
    \]
    where $X_0$ is the least significant digit (or low-order digit) and $X_{n-1}$ is the most significant digit (or high-order digit).
\end{definition}
The number system to represent the integer $x$ consists of:
\begin{citemize}
    \item The number of digits $n$
    \item A set $D_i$ of possible values for each digit $X_i$
    \item A rule of interpretation that maps between the set of digit-vectors and the set of integers.
\end{citemize}
The size of the set of integers that can be represented by a number system is at most $K$ elements, where $K$ equals:
\[
    \prod_{i = 0}^{n -1} |D_i|
\]

\begin{eg}
    Let's consider the case where $n = 6$ and $D_i = \{0,1, \ldots, 9\}$ of cardinality $|D_i| = 10$. The size of the set of integers that can be represented is at most:
    \[
        K = \prod_{i = 0}^{5} |D_i| = 10^6 = 1000000
    \]
    Thus the number system can represent at most 1 million integers, which are the integers from $0$ to $999999$.
\end{eg}

\begin{definition}[Redundance and Non-redundance]
    A number system is said to be redundant if there exist at least two distinct digit-vectors that represent the same integer. Otherwise, the number system is said to be non-redundant.
\end{definition}

\begin{definition}[Weighted (Positional) Number System]
    A weighted number system is a non-redundant number system in which each digit position has a weight (or value) associated with it. It can be denoted as:
    \[
        x = \sum_{i = 0}^{n - 1} X_i \cdot W_i
    \]
    where $W = (W_{n-1}, W_{n-2}, \ldots, W_0)$ is the weight vector of size $n$ of the number system.
\end{definition}

\begin{eg}
    Let's consider the case where $n = 6$, the digit vector $X = (8,5,4,7,0,3)$ and the weight vector $W = (10^5, 10^4, 10^3, 10^2, 10^1, 10^0)$. Thus we have:
    \[
        x = 8 \cdot 10^5 + 5 \cdot 10^4 + 4 \cdot 10^3 + 7 \cdot 10^2 + 0 \cdot 10^1 + 3 \cdot 10^0 = 854703_{10}
    \]
\end{eg}
Note that when the weight vector is of the form:
\begin{citemize}
    \item $W_0 = 1$
    \item $W_i = W_{i-1} R_{i-1}$, for $1 \leq i \leq n - 1$
\end{citemize}
it is said to be a radix number system.

\begin{definition}[Radix Number System]
    A radix number system is a weighted number system in which the weight vector is defined as follows:
    \begin{citemize}
        \item $W_0 = 1$
        \item $W_i = W_{i-1} R_{i-1}$, for $1 \leq i \leq n - 1$
    \end{citemize}
    or equivalently:
    \[
        W_i = \prod_{j = 0}^{i - 1} R_j
    \]
    where $R = (R_{n-1}, R_{n-2}, \ldots, R_0)$ is the radix vector of size $n$ of the number system.
\end{definition}
Note that the radix can be fixed and thus the weight vector can be defined as:
\[
    W_i = R^i
\]
In this case, the number system is called a fixed-radix number system. Otherwise, the number system is called a mixed-radix number system.

\begin{eg}
    Let's consider the weight vector in a fixed-radix number system with $r = 10$. We have:
    \[
        W = (10^{n-1}, 10^{n-2}, \ldots, 10, 1)
    \]
    Thus:
    \[
        x = \sum_{i = 0}^{n - 1} X_i \cdot 10^i
    \]
    which is the familiar decimal number system.
\end{eg}
Note that almost all common number systems are fixed-radix number systems, such as the binary number system (with $r = 2$), the octal number system (with $r = 8$) and the hexadecimal number system (with $r = 16$).

\begin{eg}
    Let's consider a mixed-radix number system, like the time system. The radix vector is $R = (24, 60, 60)$. Thus the weight vector is:
    \[
        W = (3600, 60, 1)
    \]
\end{eg}

\begin{definition}[Canonical Number System]
    A canonical number system is a radix number system in which the digit set $D_i$ for each digit position $i$ is defined as:
    \[
        D_i = \{0, 1, \ldots, R_i - 1\}
    \]
    where $|D_i| = R_i$ the corresponding radix of the digit position $i$.
\end{definition}

\begin{eg}
    In a $n$ fixed-radix-$r$ number system, the digit set for each digit position is:
    \[
        D_i = \{0, 1, \ldots, r - 1\}
    \]
    and thus the range of integers that can be represented is from $0$ to $r^n - 1$.
\end{eg}

\begin{definition}[Conventional Number System]
    A conventional number system is a canonical number system in which the radix is fixed and positive.
\end{definition}
Note that these are by far the most commonly used number systems.

\subsection{Transformation of Nonnegative Numbers to and from Decimal}
The conversion table (up to $15$) for the most commonly used number systems is given below:
\vskip0.3cm

\begin{center}
\begin{tabular}{
>{\centering\arraybackslash}p{0.22\textwidth} |
>{\centering\arraybackslash}p{0.22\textwidth} |
>{\centering\arraybackslash}p{0.22\textwidth} |
>{\centering\arraybackslash}p{0.22\textwidth}
}
&&\\ \textbf{Decimal} & \textbf{Binary} & \textbf{Octal} & \textbf{Hexadecimal} \\ &&\\ \hline &&\\
$0$  & $0000$ & $00$ & $0$ \\
$1$  & $0001$ & $01$ & $1$ \\
$2$  & $0010$ & $02$ & $2$ \\
$3$  & $0011$ & $03$ & $3$ \\
$4$  & $0100$ & $04$ & $4$ \\
$5$  & $0101$ & $05$ & $5$ \\
$6$  & $0110$ & $06$ & $6$ \\
$7$  & $0111$ & $07$ & $7$ \\
$8$  & $1000$ & $10$ & $8$ \\
$9$  & $1001$ & $11$ & $9$ \\
$10$ & $1010$ & $12$ & $A$ \\
$11$ & $1011$ & $13$ & $B$ \\
$12$ & $1100$ & $14$ & $C$ \\
$13$ & $1101$ & $15$ & $D$ \\
$14$ & $1110$ & $16$ & $E$ \\
$15$ & $1111$ & $17$ & $F$ \\
\end{tabular}
\end{center}
Note that generally in programming languages the prefix $0\text{x}$ is used to denote hexadecimal numbers.

\begin{eg}
    Let's convert $10011_2$ to decimal. We have:
    \[
        10011_2 = 1 \cdot 2^4 + 0 \cdot 2^3 + 0 \cdot 2^2 + 1 \cdot 2^1 + 1 \cdot 2^0 = 19_{10}
    \]
    Let's now convert $179_{10}$ to binary. We have:
    \begin{align*}
        179 \div 2 &= 89 \text{ remainder } 1 \\
        89 \div 2 &= 44 \text{ remainder } 1 \\
        44 \div 2 &= 22 \text{ remainder } 0 \\
        22 \div 2 &= 11 \text{ remainder } 0 \\
        11 \div 2 &= 5 \text{ remainder } 1 \\
        5 \div 2 &= 2 \text{ remainder } 1 \\
        2 \div 2 &= 1 \text{ remainder } 0 \\
        1 \div 2 &= 0 \text{ remainder } 1
    \end{align*}
    Thus we have:
    \[
        179_{10} = 10110011_2
    \]
\end{eg}

\begin{eg}
    Similarly, we can convert $751_{10}$ to octal. We have:
    \begin{align*}
        751 \div 8 &= 93 \text{ remainder } 7 \\
        93 \div 8 &= 11 \text{ remainder } 5 \\
        11 \div 8 &= 1 \text{ remainder } 3 \\
        1 \div 8 &= 0 \text{ remainder } 1
    \end{align*}
    Thus we have:
    \[
        751_{10} = 1357_8
    \]
\end{eg}

\subsection{Octal or Hexadecimal to and from Binary}
\begin{definition}[Bit-Vector Representation]
    A bit-vector representation of a number is a binary representation of the number in which each digit of the original number is represented by a fixed number of bits (binary digits).
\end{definition}
Note that for powers of $2$ radix, a digit $d$ is represented by a bit-vector $(d_{k -1}, \ldots, d_0)$ where $k = \log_2 r$ bits such that:
\[
    d = \sum_{i = 0}^{k - 1} d_i \cdot 2^i
\]
For example, the digit $D_{16}$ in the hexadecimal number system is represented by a $4$-bit binary vector $1101_2$.

\begin{eg}
    Let's convert $10000100110_2$ to octal. Since $k = \log_2 8 = 3$, we group the bits into groups of $3$ starting from the right:
    \[
        10000100110_2 = 10 \ 000 \ 100 \ 110_2 = 010 \ 000 \ 100 \ 110_2 = 2046_8
    \]
    Note that we added leading zeros to the leftmost group to make it a full group of $3$ bits.
\end{eg}

\begin{eg}
    Let's convert $2543_8$ to binary. Since $k = \log_2 8 = 3$, each octal digit is represented by a $3$-bit binary vector:
    \[
        2543_8 = 010 \ 101 \ 100 \ 011_2 = 101011100011_2
    \]
\end{eg}
The same process can be applied to convert between hexadecimal and binary, but with groups of $4$ bits instead of $3$ bits.

\section{Representation of Signed Integers}
\begin{definition}[Sign-and-Magnitude (SM) Representation]
    A signed integer is represented by a pair:
    \[
        (x_s,x_m)
    \]
    where $x_s$ is the sign bit and $x_m$ is the magnitude of the integer. The sign bit is typically defined as:
    \[
        x_s = \begin{cases}
        0 & \text{if } x \geq 0 \\
        1 & \text{if } x < 0
        \end{cases}
    \]
    and the magnitude can be represented as any nonnegative integer.
\end{definition}
Note that in this representation, the integer $0$ can be represented in two ways: $(0,0)$ and $(1,0)$, which makes it a redundant representation.

\begin{eg}
    Traditionally, in an $n$-bit sign-and-magnitude representation, the most significant bit (MSB) is used as the sign bit and the remaining $n - 1$ bits are used to represent the magnitude. For example:
    \begin{align*}
        5_{10} &= (0, 101)_2 = 0101_2 \\
        -5_{10} &= (1, 101)_2 = 1101_2 \\
        0_{10} &= (0, 000)_2 \text{ or } (1, 000)_2
    \end{align*}
\end{eg}
Such a representation have some advantages such as being symmetrical around zero, the magnitude lies between $-(2^{n-1} - 1)$ and $2^{n-1} - 1$, and being easy to understand. However, it has some disadvantages such as having complex arithmetic operations and having two representations for zero.

\subsection{True-and-Complement Representation}
\begin{definition}[True-and-Complement (TC) Representation]
    A signed integer is represented by a pair:
    \[
        (x_s,x_m)
    \]
    where $x_s$ is the sign bit and $x_m$ is the magnitude of the integer. The sign bit is typically defined as:
    \[
        x_s = \begin{cases}
        0 & \text{if } x \geq 0 \\
        1 & \text{if } x < 0
        \end{cases}
    \]
    and the magnitude can be represented as any nonnegative integer. However, in this representation, the magnitude of negative integers is represented by the complement of their absolute value.
\end{definition}